
\hzsync{equation}{73}
\hzsync{theorem}{8}
\hzsync{proposition}{7}

\section{\trPositiveDefinite\trQuadraticForm{}的\trSymplectic{}分类}
\label{sec:quadratic-forms-classification}

在转向借助\trSymplecticMapping{}把一个\trOpenSet{}\trEmbedding{}于另一个%
\trOpenSet{}之中这一微妙问题之前，先考察如下特殊情形将是有益的：所述%
\trOpenSet{}为\textbf{\trEllipsoid}（ellipsoid），而所考虑的\trSymplecticMapping{}%
是\trLinear{}的。我们从一个在力学中人所共知的代数考察入手。回忆：向量空间 $V$%
上的函数 $q:V\to\R$ 称为\textbf{\trQuadraticForm}（quadratic form），如果存在%
\trSymmetric{}的\trBilinear{}形式 $\hat{q}:V\times V\to\R$ 使得
\[
q(x)=\frac{1}{2}\,\hat{q}(x,x),\qquad x\in V.
\]
\trQuadraticForm{} $q$ 称为\trPositiveDefinite{}的，如果 $\hat{q}$ 是 $V$ 中的%
一个\trInnerProduct。在 $V$ 中取一组基，便有表示
\begin{align*}
\hat{q}(x,y) &= \langle Sx,\ y\rangle\\
q(x) &= \frac{1}{2}\langle Sx,\ x\rangle,\qquad S=S^{T}>0
\end{align*}
其中 $x$, $y\in V$。从现在起，$(V,\omega)=(\R^{2n},\omega)$ 是%
\trStandard\trSymplecticVectorSpace，其\trSymplecticForm{}（此前记作 $\omega_0$）%
由
\[
\omega(x,y)=\langle Jx,\ y\rangle,\qquad x,\ y\in\R^{2n}
\]
给出，而
\[
J=\begin{pmatrix}0 & 1\\ -1 & 0\end{pmatrix}.
\]

把 $V$ 上\trPositiveDefinite\trQuadraticForm{}的全体记作 $\mathcal{P}$，并把 $V$%
的\trLinear\trSymplectic\trMapping{} $\varphi$（即满足 $\varphi^{*}\omega=\omega$ 者）%
所成的群记作 $G=\Sp(n)$，那么我们有由复合
\[
\mathcal{P}\times G\to\mathcal{P}\ :\quad[q,\varphi]\mapsto q\circ\varphi.
\]
所定义的 $G$ 在 $\mathcal{P}$ 上的作用。结果表明，$\mathcal{P}$ 在此作用下的%
\trOrbit{}由 $n$ 个正数刻画。借助于一个使形式 $\omega$ 保持\trInvariant{}的%
\trLinear\trMapping{}，形式 $q$ 可以\trDiagonal{}化。这一代数事实早已为 K.~\trWeierstrass{} [220]%
所知。

\begin{theorem}\label{thm:quadratic-normal-form}
若 $q\in\mathcal{P}$，则存在\trLinear\trMapping{} $\varphi$ 满足 $\varphi^{*}\omega=\omega$，%
使得
\[
q\circ\varphi(x)=\frac{1}{2}\sum_{j=1}^{n}\lambda_j\bigl(x_j^{2}+x_{n+j}^{2}\bigr),
\]
其中 $0<\lambda_n\leq\lambda_{n-1}\leq\dots\leq\lambda_1$。诸数 $\lambda_j$ 由 $q$%
与 $\omega$ 唯一确定。事实%
上，$(\pm i\lambda_j)$（$1\leq j\leq n$）恰是\trLinear\trHamiltonianVectorField{}%
$X_q$ 的\textbf{\trEigenvalue}（eigenvalue），此 $X_q$ 由关系式
\[
\omega\bigl(X_q(x),\,y\bigr)=-dq(x)(y)
\]
（对一切 $x$, $y\in\R^{2n}$）与 $q$ 及 $\omega$ 相联系。
\end{theorem}

我们把证明稍作搁置，先来指出：若 $q(x)=\frac{1}{2}\langle Sx,\,x\rangle$，其中 $S$%
为\trPositiveDefinite\trSymmetric{}矩阵，则 $X_q$ 是\trLinear\trVectorField{}%
$X_q(x)=JSx$；于是，若 $\tau\in\Sp(n)$ 是任一\trSymplecticMatrix，则由%
\trHamiltonianVectorField{}的变换公式即得
\[
X_{q\circ\tau}=\tau^{-1}\cdot X_q\cdot\tau.
\]
于是，这两个\trLinear\trVectorField{} $X_q$ 与 $X_{q\circ\tau}$ 彼此相似，从而具有%
相同的\trEigenvalue。把该定理所保证的 $n$ 个正实数 $\lambda=\lambda(q)\in\R^{n}$%
与 $q$ 相联系，我们便得出
\[
\lambda(q\circ\tau)=\lambda(q)
\]
对一切 $q\in\mathcal{P}$ 与 $\tau\in\Sp(n)$ 成立，因此这些数刻画了 $\mathcal{P}$%
在 $\Sp(n)$ 作用下的\trOrbit。

这当然与满足 $A^{T}A=1$ 的矩阵所组成的\trOrthogonal{}群（orthogonal group）%
的作用形成对照。由于 $q\circ A(x)=\frac{1}{2}\langle SAx,\ Ax\rangle=%
\frac{1}{2}\langle A^{T}SAx,\ x\rangle=\frac{1}{2}\langle A^{-1}SAx,\ x\rangle$，%
我们看到，$\mathcal{P}$ 在 $O(2n)$ 作用下的\trOrbit{}由\trSymmetric{}矩阵 $S$%
的 $2n$ 个正\trEigenvalue{}决定，这与 $JS$ 的 $n$ 个虚\trEigenvalue{}偶对%
$(\pm i\lambda_j)$ 恰成对照。例如在 $\R^{2}$ 上，一个正\trQuadraticForm{}经过%
\trOrthogonal{}变换后形如 $q(x)=\frac{1}{2}\langle Sx,\,x\rangle$，其中%
$S=\operatorname{Diag}(\lambda_1,\lambda_2)$ 为\trDiagonal{}矩阵%
且 $\lambda_j>0$。该定理保证存在满足 $\det\tau=1$ 的\trLinear\trMapping{} $\tau$，%
把 $q$ 变为\trSymplectic\trNormalForm{} $q\circ\tau(x)=\frac{1}{2}\langle Qx,\,x\rangle$，%
其中 $Q=\operatorname{Diag}(\lambda,\lambda)$。几何上，\trEllipsoid{}%
$\{q(x)<1\}$ 被\trSymplectic{}地映到具有相同面积的\textbf{\trDisc}（disc）%
$\{q\circ\tau(x)<1\}$ 上。我们看到，在 $\R^{2}$ 中诸\trNormalForm{}对应于%
由面积区分的\trDisc。

\begin{proof}
在\textbf{\trProductSpace}（product space） $(x,y)\in V\times V$ 上，函数 $K(x,y)=q(x)+q(y)$ 与 $\omega(x,y)$%
都是\trQuadraticForm。由于 $K$ 是\trPositiveDefinite{}的，我们可以应用一个%
人所共知的\trVariational{}论证。在由
\[
M=\bigl\{(x,y)\in V\times V\ \bigm|\ q(x)+q(y)=1\bigr\}
\]
所定义的\trCompact\trManifold{} $M=M(q)$ 上，连续函数%
$(x,y)\mapsto\omega(x,y)$ 达到最大值，设在 $(a,b)\in M$ 处取得，使得
\[
\omega(a,b)=\max_{(x,y)\in M}\omega(x,y).
\]
我们指出，$\omega|_{M}$ 的\textbf{\trCriticalPoint}（critical point）$(a,b)\in M$%
位于一整个由属于同一\trCritical{}水平（critical level）的\trCriticalPoint{}%
所组成的\trCircle{}之上。事实%
上，若把 $V\times V$ 上的\trCircle{}作用（circle action）定义为
\[
e^{it}\cdot(x,y)=\bigl(x(t),\,y(t)\bigr),\qquad t\in\R,
\]
其中
\[
x(t)=(\cos t)x+(\sin t)y,\qquad y(t)=(\cos t)y-(\sin t)x,
\]
则容易验证：诸\trQuadraticForm{}在该作用下保持\trInvariant，即
\[
K\bigl(x(t),\,y(t)\bigr)=K(x,y)\qquad\text{且}\qquad
\omega\bigl(x(t),\,y(t)\bigr)=\omega(x,y)
\]
对一切 $(x,y)\in V\times V$ 成立。

若 $\lambda\in\R$ 表示\trCriticalPoint{} $(a,b)\in V\times V$ 的\trLagrange{}乘子%
（Lagrange multiplier），则诸\trDerivative{}满足%
$dK(a,b)=\lambda\,d\omega(a,b)$，明显写出即为
\begin{equation}\label{eq:lagrange-multiplier}
\begin{aligned}
\hat{q}(a,x) &= dq(a)(x) = \lambda\,\omega(x,b)\\
\hat{q}(y,b) &= dq(b)(y) = \lambda\,\omega(a,y)
\end{aligned}
\end{equation}
其中 $x$, $y\in V$ 任意。这里我们用到了 $q$ 的\trDerivative{}的表示%
$dq(x)(y)=\hat{q}(x,y)$。鉴于 $q(x)=\frac{1}{2}\hat{q}(x,x)$，从%
\eqref{eq:lagrange-multiplier} 可得出%
$2=2\bigl(q(a)+q(b)\bigr)=2\lambda\,\omega(a,b)$。由于最大值 $\omega(a,b)$ 为正，
\begin{equation}\label{eq:omega-ab-positive}
\omega(a,b)=\frac{1}{\lambda}>0.
\end{equation}
特别地，由于 $\omega$ 是\trAntisymmetric{}的，$a$, $b$ 在 $V$ 中%
\trLinearlyIndependent。由于 $\omega(a,b)\neq0$，二维\trSubspace{}%
$V_1=\operatorname{span}\{a,b\}\subset V$ 是\trSymplectic{}的。它是 $X_q$%
的属于纯虚\trEigenvalue{}对 $\pm i\lambda$ 的\textbf{\trEigenspace}（eigenspace）。事实%
上，方程 \eqref{eq:lagrange-multiplier} 等价于
\[
X_q(a)=+\lambda b\qquad\text{且}\qquad X_q(b)=-\lambda a,
\]
从而 $X_q(a\pm ib)=\mp(i\lambda)(a\pm ib)$。经规范化，我们引入%
$e_1=\alpha a$, $f_1=-\alpha b$（$\alpha^{2}=\lambda$），并求得%
$\omega(f_1,e_1)=1$。于是，偶对 $e_1$, $f_1$ 构成 $V_1$ 的一组%
\trSymplecticBasis，对它而言
\begin{equation}\label{eq:q-e1-f1}
q(e_1)=q(f_1)=\frac{\lambda}{2}\qquad\text{且}\qquad\hat{q}(e_1,f_1)=0.
\end{equation}
回忆对 $x$, $y\in V$ 的恒等式 $\hat{q}(x,y)=q(x+y)-q(x)-q(y)$，我们便有
\begin{equation}\label{eq:q-normal-v1}
q(v)=\frac{\lambda}{2}\bigl(\xi^{2}+\eta^{2}\bigr)
\end{equation}
其中 $v=\xi e_1+\eta f_1\in V_1$；这正是所要的 $V_1$ 中的\trNormalForm。%
为了迭代这一论证，我们把
\[
V_1^{\perp}=\bigl\{\,v\in V\ \bigm|\ \omega(v_1,\,v)=0
\quad\text{对一切}\quad v_1\in V_1\,\bigr\}
\]
称为 $V$ 中的\textbf{\trSymplecticComplement}（symplectic complement）。于是 $V_1^{\perp}$%
是一个\trSymplectic\trSubspace，并且
\[
V=V_1\oplus V_1^{\perp}.
\]
事实%
上，$V_1\cap V_1^{\perp}=\{0\}$，因为 $\omega$ 在 $V_1$ 上\trNondegenerate；%
此外，对每个 $v\in V$ 都有%
$v-\omega(v,e_1)f_1+\omega(v,f_1)e_1\in V_1^{\perp}$，从而%
$V=V_1+V_1^{\perp}$。若 $v\in V_1$ 且 $w\in V_1^{\perp}$，则 $\omega(v,w)=0$，%
而且我们断言：就 $\hat{q}$ 而言 $v$ 与 $w$ 也是\trOrthogonal{}的：
\[
\hat{q}(v,w)=0.
\]
确实如此，因为 $V_1$ 是一个\trEigenspace。事实%
上，把 $x$ 与 $y\in V_1^{\perp}$ 代入方程组 \eqref{eq:lagrange-multiplier}，%
右端便化为零，从而左端亦化为零；又因 $V_1=\operatorname{span}\{a,b\}$，%
断言得证。等价地，我们有
\begin{equation}\label{eq:q-additive}
q(v+w)=q(v)+q(w)
\end{equation}
其中 $v\in V_1$ 且 $w\in V_1^{\perp}$。

我们可以把与上面相同的论证应用于 $q$ 在\trSymplectic{}空间 $V_1^{\perp}$%
上的限制。迭代地，我们便求得一个分解%
$V=V_1\oplus V_2\oplus\dots\oplus V_n$，它把 $V$ 分成 $n$ 个特定的二维%
\trSymplectic{}空间 $V_j$；这些空间是 $X_q$ 的属于\trEigenvalue{}对%
$(\mp i\lambda_j)$ 的\trEigenspace，其中 $\lambda_j>0$。鉴于%
\eqref{eq:q-additive}，
\[
q(v)=q(v_1)+q(v_2)+\dots+q(v_n)
\]
其中 $v=v_1+v_2+\dots+v_n$，$v_j\in V_j$。此外，在 $V_j$ 中存在%
\trSymplecticBasis{} $(e_j,f_j)$，使得
\[
q(v)=\frac{1}{2}\sum_{j=1}^{n}\lambda_j\bigl(x_j^{2}+x_{n+j}^{2}\bigr)
\]
其中 $v=\sum_{j=1}^{n}\bigl(x_j e_j+x_{n+j}f_j\bigr)$。这就是 $q$ 所要的%
\trNormalForm，至此\cref{thm:quadratic-normal-form}的证明即告完成。
\end{proof}

改换记号并引入 $\lambda_i=\tfrac{2}{r_i^{2}}$（$i=1,2,\dots,n$），我们把该%
\trNormalForm{}写作
\[
q\circ\varphi(x)=\sum_{j=1}^{n}\frac{1}{r_j^{2}}\bigl(x_j^{2}+x_{j+n}^{2}\bigr),
\]
从而给每个 $q\in\mathcal{P}$ 配上了正数
\[
0<r_1(q)\leq r_2(q)\leq\dots\leq r_n(q),
\]
它们刻画了 $\Sp(n)$ 在 $\mathcal{P}$ 上的作用之\trOrbit。简记%
$r=(r_1,\,r_2,\,\dots,\,r_n)$，我们已经证明了：

\begin{proposition}\label{prop:r-composition-invariant}
$r(q)=r(q\circ\varphi)$ 对一切 $q\in\mathcal{P}$ 与一切 $\varphi\in\Sp(n)$ 成立。%
此外，若 $r(q)=r(Q)$，其中 $q$, $Q\in\mathcal{P}$，则存在满足 $Q=q\circ\varphi$%
的 $\varphi\in\Sp(n)$。
\end{proposition}

\begin{proposition}\label{prop:quadratic-monotonicity}
设 $q$ 与 $Q$ 属于 $\mathcal{P}$。那么
\[
Q\leq q\quad\Longrightarrow\quad r_j(q)\leq r_j(Q)
\qquad\text{对}\quad 1\leq j\leq n.
\]
\end{proposition}
\begin{proof}
从\cref{thm:quadratic-normal-form}的证明出发并回忆上面的记号，我们知道
\[
\frac{1}{2}\,r_n(q)^{2}=\frac{1}{\lambda_n(q)}=\max_{M(q)}\omega=\omega(a,b),
\]
其中 $M(q)=\bigl\{(x,y)\in V\times V\ \bigm|\ q(x)+q(y)=1\bigr\}$。由于对一切%
$x\in V$ 有 $Q(x)\leq q(x)$，故存在 $\alpha\geq1$ 使得%
$1=Q(\alpha a)+Q(\alpha b)$。于是%
$(\alpha a,\alpha b)\in M(Q)=\bigl\{(x,y)\ \bigm|\ Q(x)+Q(y)=1\bigr\}$，%
并且 $\omega(\alpha a,\alpha b)=\alpha^{2}\,\omega(a,b)\geq\omega(a,b)$。%
取 $\omega$ 在 $M(Q)$ 上的最大值，我们便得出 $r_n(Q)\geq r_n(q)$。更一般地，%
利用\trCritical{}水平的下述\trMinimax{}刻画，容易验证\trEstimate{}%
$r_j(Q)\geq r_j(q)$（$1\leq j\leq n$）：
\[
\frac{1}{2}\,r_j(q)^{2}=\min_{W\in F_j}\Bigl\{\ \max_{W\cap M(q)}\omega\ \Bigr\},
\]
其中 $F_j$ 表示由维数为 $2n+2j$ 的一切\trLinear\trSubspace{}%
$W\subset V\times V$ 所组成的族，而 $1\leq j\leq n$。事实%
上，考虑到 $\omega$ 与 $K$ 在\trProductSpace{} $V\times V$ 上都是\trQuadraticForm{}%
且 $K$ 是\trPositiveDefinite{}的，这不过是\trCourant--\trHilbert{}\trMinimax{}原理，%
即关于\trSymmetric{}矩阵之\trEigenvalue{}的原理，该矩阵相对于由正形式 $K$%
所定义的\trInnerProduct{}与形式 $\omega$ 相联系。然而，我们顾及到了该问题在%
上文所定义的\trCircle{}作用之下是\trInvariant{}的。特别地，若 $(a,b)\in V\times V$ 是一个%
\trCriticalPoint，则 $(-b,a)$ 也是同一\trCritical{}水平上的\trCriticalPoint。%
于是，\trCourant--\trHilbert{}原理中的每个\trEigenvalue{}的\textbf{\trMultiplicity}（multiplicity）%
至少为 $2$。这个一般原理将给出诸正的\trCharacteristic{}水平%
（characteristic level）$c_{4n}\geq c_{4n-1}\geq\dots\geq c_{2n+1}$，%
而为了找出全部正的\trCritical{}水平，我们已在上文中选取了
\[
\frac{1}{2}\,r_j(q)^{2}=c_{2n+2j}
\]
（$1\leq j\leq n$）。
\end{proof}

为了从几何上解释上述结果，我们给形式 $q\in\mathcal{P}$ 配上\trEllipsoid{}%
$E(q)$，把它定义为\trOpenSet
\begin{equation}\label{eq:ellipsoid-definition}
E(q)=\bigl\{\,x\ \bigm|\ q(x)<1\,\bigr\}\subset\R^{2n}.
\end{equation}
于是，$\Sp(n)$ 的\textbf{\trGroupAction}（group action）便可以直观地显示为
\begin{equation}\label{eq:ellipsoid-action}
\varphi^{-1}\bigl(E(q)\bigr)=E(q\circ\varphi).
\end{equation}
注意到
\begin{equation}\label{eq:ellipsoid-inclusion}
E(q)\subset E(Q)\quad\Longleftrightarrow\quad q\geq Q,
\end{equation}
因此 $E(q)=E(Q)$ 当且仅当 $q=Q$。这样，我们便可以给每个\trOpenSet{} $E$%
配上数 $r(E)$，其定义为
\begin{equation}\label{eq:r-of-ellipsoid}
r(E)=r(q)\qquad\text{若}\quad E=E(q).
\end{equation}
鉴于 \eqref{eq:ellipsoid-action}，我们从\cref{prop:r-composition-invariant}%
得出：诸数 $r$ 在\trLinear\trSymplectic\trMapping{}所成群的作用下\trInvariant。%
概括起来，我们已经证明了：

\begin{proposition}\label{prop:ellipsoid-invariants}
若 $E$, $F$ 是\trEllipsoid{}且 $\varphi\in\Sp(n)$，则%
$r\bigl(\varphi(E)\bigr)=r(E)$。反之，若 $r(E)=r(F)$，则存在满足%
$\varphi(E)=F$ 的\trMapping{} $\varphi\in\Sp(n)$。
\end{proposition}

现在的关键观察是：这些\trInvariant{}量具有\trMonotonicity。事实%
上，鉴于 \eqref{eq:ellipsoid-inclusion}，\cref{prop:quadratic-monotonicity}%
可以重述如下。

\begin{proposition}\label{prop:ellipsoid-monotonicity}
若 $E$, $F$ 是\trEllipsoid，则
\[
E\subset F\quad\Longrightarrow\quad r_j(E)\leq r_j(F)
\qquad\text{对}\quad 1\leq j\leq n.
\]
反之，若 $r_j(E)\leq r_j(F)$（$1\leq j\leq n$），则存在满足%
$\varphi(E)\subset F$ 的 $\varphi\in\Sp(n)$。
\end{proposition}

总括我们的考察：我们对 $\R^{2n}$ 中\trEllipsoid{}这类\trOpenSet{} $E$ 定义了%
\trInvariant{}量 $r(E)$，借助它们，\trEllipsoid{}的\trLinear%
\trSymplecticEmbedding{}问题便可回答如下：

\begin{theorem}\label{thm:ellipsoid-linear-embedding}
设 $E$ 与 $F$ 是两个\trEllipsoid。则存在 $\varphi\in\Sp(n)$ 使得
\[
\varphi(E)\subset F
\]
当且仅当 $r_j(E)\leq r_j(F)$ 对一切 $1\leq j\leq n$ 成立。
\end{theorem}

我们用一个例子来阐释这一论断。设%
$B=B(R)=\bigl\{x\ \bigm|\ |x|^{2}<R^{2}\bigr\}$ 是半径为 $R$ 的开\trBall，%
则 $r_j(B)=R$ 对一切 $1\leq j\leq n$ 成立，而且当且仅当 $E$ 受条件%
$r_1(E)\geq R$ 限制时，我们才能找到把 $B$ \trEmbedding{}到\trEllipsoid{} $E$%
之内的\trMapping{} $\varphi\in\Sp(n)$。这一切当然不足为奇。毕竟我们考虑的只是%
\trLinear\trSymplectic\trMapping{}，人们或许会指望，使用\trNonlinear%
\trSymplecticEmbedding{}会更加灵活。然而结果表明——这是一个直到晚近才被%
发现的、十分令人惊讶的现象——只要 $\varphi:B(R)\to E$ 是任意一个%
\trSymplecticEmbedding，就仍有 $r_1(E)\geq R$。稍后我们将利用一个整体%
\trSymplecticInvariant{} $c(U)$ 来证明这一\trRigidity{}论断：这个 $c(U)$%
对每个\trOpenSet{} $U\subset\R^{2n}$ 都有定义，并且在\trSymplecticDiffeomorphism{}%
之下满足 $c\bigl(\varphi(U)\bigr)=c(U)$。这一\trInvariant{}量扩充了 $r_1$，%
使得对\trEllipsoid{} $E$ 有 $c(E)=\pi r_1(E)^{2}$；此外，它还具有下述意义下的%
\trMonotonicity：只要存在\trSymplecticEmbedding{} $\varphi:U\to V$，%
便有 $c(U)\leq c(V)$。

由于这一整体\trInvariant{}量 $c$ 的构造将与\trHamiltonianEquation{}的特殊%
\trPeriodicSolution{}相关联，从动力学的观点来诠释\trSymplecticInvariant{}%
$r(E)$ 将是有益的。$r(E)$ 与 $\partial E$ 上某些\trClosedCharacteristic{}%
之间存在着一种关系；这里 $\partial E$ 是 $E$ 的\trBoundary，它是一个%
\trCompact\trHypersurface：
\[
\partial E=\bigl\{\,z\ \bigm|\ q(z)=1\,\bigr\},\qquad\text{若}\ E=E(q).
\]
我们可以假定 $q$ 已呈\trNormalForm：
\[
q(z)=\sum_{j=1}^{n}\frac{x_j^{2}+y_j^{2}}{r_j^{2}}.
\]
相应地，$V$ 就分裂成二维\trSymplectic{}平面%
$V=V_1\oplus V_2\oplus\dots\oplus V_n$，而 $z_j=(x_j,y_j)\in\R^{2}$ 是 $V_j$%
中的\trSymplecticCoordinates。与之相联系的\trHamiltonianVectorField{} $X_q$%
的解 $z(t)=\bigl(z_1(t),\,\dots,\,z_n(t)\bigr)$ 容易算出。令%
$\lambda_j=\tfrac{2}{r_j^{2}}$，便有
\[
z_j(t)=e^{t\lambda_j J}z_j(0),\qquad 1\leq j\leq n,
\]
其中 $J$ 是\trSymplecticStructure{}的\trStandard{} $2\times 2$ 矩阵。于是，%
$\partial E\cap V_j$ 上带有一条特定的\trPeriodicSolution{}%
$w_j(t)=\bigl(0,\dots,0,\,z_j(t),\,0,\dots,0\bigr)$，它由
\[
z_j(t)=e^{t\lambda_j J}z_j(0)\qquad\text{与}\qquad
\bigl|z_j(0)\bigr|^{2}=r_j^{2}
\]
给出，并且其\trPeriodic{}为 $T=\pi r_j^{2}$。它的\trAction{}——一个整体%
\trSymplecticInvariant{}——便是
\[
A(w_j)=\frac{1}{2}\int_{0}^{T}\bigl\langle-J\dot{w}_j,\,w_j\bigr\rangle
=\pi r_j^{2}.
\]
我们看到，集合 $E$ 的\trInvariant{}量 $r_j$ 由 $\partial E$ 上的一条%
\trClosedCharacteristic{}表示出来。特别地，
\[
\pi r_1(E)^{2}=\min\ \Bigl\{\ \bigl|A(z)\bigr|\ :\ z\ \text{为}\
\partial E\ \text{上的\trClosedCharacteristic}\ \Bigr\}.
\]
在下文中，这一动力学方法对于构造\trSymplecticInvariant{}将极其有用。

我们还想补充一点：如果人们在一个\textbf{\trEquilibrium}（equilibrium）附近局部地研究%
\trHamiltonianVectorField{} $X_H$ 的\trOrbit{}结构，\trQuadraticForm{}便起着%
重要的作用；所谓\trEquilibrium，即 $X_H$ 在其上为零的点。我们将通过%
描述阐释这一点的若干众所周知的结果来结束本章。我们可以假定\trEquilibrium{}%
为原点，于是 $\nabla H(0)=0$，并且
\[
X_H(z)=JH_{zz}(0)z+\dotsb
\]
对 $\R^{2n}$ 中原点附近的 $z$ 成立。从现在起，我们将假定 $JH_{zz}(0)$ 的%
\trEigenvalue{}都是纯虚的，于是 $(\pm i\lambda_j)$（$1\leq j\leq n$）便是%
全部\trEigenvalue。如果进而假定 $JH_{zz}(0)$ 可以\trSymplectic{}地\trDiagonal{}化，%
那么经过一个\trLinear\trSymplectic{}变换便可写出
\[
H(z)=\frac{1}{2}\sum_{j=1}^{n}\lambda_j\bigl(x_j^{2}+y_j^{2}\bigr)+\dotsb,
\]
其中 $\lambda_j\in\R$ 为实数。鉴于\cref{thm:quadratic-normal-form}，只要%
$H_{zz}(0)$ 是\trPositiveDefinite{}的，我们就总能达到这一\trNormalForm。%
这时，$\R^{2n}$ 分解为一些\trSymplectic{}二维平面 $(x_j,y_j)$ 之和，%
这些平面充满了\trLinear{}系统 $\dot{z}=JH_{zz}(0)z$ 的\trPeriodicSolution，%
其\trPeriodic{}为 $\tfrac{2\pi}{|\lambda_j|}$。随之而来的问题是：这些解族%
是否可以\textbf{\trContinuation}（continuation）为 $z=0$ 附近\trNonlinear{}%
系统的解？正如例子所表明的，\trVectorField{} $X_H$ 可能不容许除 $z=0$%
之外的任何\trPeriodicSolution。通过对诸 $\lambda_j$ 加以限制，一个归功于%
\trLyapunov{}[145] 的众所周知的结果保证了在无穷多条%
\textbf{\trNonresonance{}条件}（nonresonance conditions）之下的%
\trContinuation。因此，十分不寻常的是：在上面研究过的\trPositiveDefinite{}%
情形中，根本不需要任何这类条件，下述归功于 A.~\trWeinstein{}[223] 的结果%
正表明了这一点：

\begin{theorem}\label{thm:weinstein-local-existence}
设 $H_{zz}(0)>0$。则对充分小的 $\varepsilon>0$，在\trHypersurface{}%
$H(z)-H(0)=\varepsilon^{2}$ 上至少有 $n$ 条\trClosedCharacteristic。作为系统%
$X_H$ 的\trParametrization{}解，它们的\trPeriodic{}接近\textbf{\trLinearization}（linearization）系统%
$\dot{z}=JH_{zz}(0)z$ 的\trPeriodic。
\end{theorem}

这一局部存在性论断已被 J.~\trMoser{}[163] 所推广。到目前为止所提到的%
\trPeriodicSolution{}都是由\trLinearization{}方程的解所引发的，而现在我们要转向%
真正的\trNonlinear{}现象。首先我们回顾 $H$ 在 $z=0$ 处的一种%
\trNonlinear\trNormalForm，它可追溯到 G.~\trBirkhoff{}（[26]，1927），%
并且建立起在\trEquilibrium{}处与 $H$ 相联系的\trNonlinear{}局部%
\trSymplecticInvariant。
